% !TeX root = main.tex
\newcommand{\imageplaceholder}[3]{%
  \begin{figure}[htbp]
    \centering
    \IfFileExists{images/#1}{%
      \includegraphics[width=0.9\textwidth]{#1}%
    }{%
      \fbox{%
        \parbox[c][3cm][c]{0.85\textwidth}{%
          \centering
          Место для рисунка\\[0.4em]
          \texttt{images/\detokenize{#1}}
        }%
      }%
    }
    \caption{#2}
    \label{#3}
  \end{figure}
}

\tableofcontents
\newpage

\section{Цель работы}

Освоить принципы построения адаптивной системы управления линейным
многомерным объектом по состоянию с использованием эталонной модели.

\section{Исходные данные}

Объект управления:
\begin{equation}
  \dot{\bm x}=A\bm x+\bm b u,
  \qquad y=C\bm x
\end{equation}
Для варианта 1 заданы
\begin{equation}
  A=\begin{pmatrix}0&1\\1&1\end{pmatrix},
  \qquad
  \bm b=\begin{pmatrix}0\\1\end{pmatrix},
  \qquad C=\begin{pmatrix}1&0\end{pmatrix},
  \qquad b_0=1
\end{equation}
Задающее воздействие и требуемые показатели качества:
\begin{equation}
  g(t)=3\sign\bigl(\cos(0{,}2t)\bigr)+3,
  \qquad t_{\text{п}}=0{,}16\,\text{с},
  \qquad \sigma=0\%
\end{equation}
Цель управления:
\begin{equation}
  \bm e=\bm x_M-\bm x,
  \qquad \lim_{t\to\infty}\bm e(t)=\bm 0
\end{equation}
Построим графики переходной функции:

\section{Эталонная модель}

\begin{equation}
  \dot{\bm x}_M=A_M\bm x_M+\bm b_M g,
  \qquad y_M=C\bm x_M
\end{equation}
\begin{equation}
  A_M=\begin{pmatrix}0&1\\-a_{M0}&-a_{M1}\end{pmatrix},
  \qquad \bm b_M=\begin{pmatrix}0\\a_{M0}\end{pmatrix}
\end{equation}

% Выбрать стандартный характеристический полином и рассчитать коэффициенты
% по заданным времени переходного процесса и перерегулированию.
\begin{equation}
  a_{M0}=\ldots,
  \qquad a_{M1}=\ldots
\end{equation}

% Переходная характеристика: единичное ступенчатое воздействие,
% нулевые начальные условия; на графике отметить t_п и sigma.
\imageplaceholder{reference_step.png}
  {Переходная характеристика эталонной модели}
  {fig:reference-step}

\noindent\textbf{Вывод по пункту:} \ldots

\section{Неадаптивная система управления}

Условия согласования и закон управления при известных параметрах объекта:
\begin{equation}
  A_M=A+\bm b\bm\theta^{\mathsf T},
  \qquad \bm b=\kappa\bm b_M,
  \qquad u=\bm\theta^{\mathsf T}\bm x+\frac{1}{\kappa}g
\end{equation}

% Рассчитать номинальные коэффициенты регулятора и проверить устойчивость.
\begin{equation}
  \bm\theta_{\text{ном}}=\ldots,
  \qquad \kappa=\ldots
\end{equation}
Во всех трёх опытах коэффициенты регулятора сохраняют номинальные значения.
Для каждого опыта приводятся графики $\bm x(t)$ и $\bm x_M(t)$
на одном поле и $\bm e(t)$ на другом.

\subsection{Номинальные параметры объекта}

\begin{equation}
  A^{(1)}=A
\end{equation}
\imageplaceholder{nonadaptive_nominal.png}
  {Неадаптивная система при номинальных параметрах объекта}
  {fig:nonadaptive-nominal}

\subsection{Отклонение параметров с сохранением устойчивости}

\begin{equation}
  A^{(2)}=\begin{pmatrix}0&1\\\ldots&\ldots\end{pmatrix}
\end{equation}
% Выбрать небольшое отклонение параметров и проверить устойчивость
% матрицы A^(2) + b theta_ном^T.
\imageplaceholder{nonadaptive_stable.png}
  {Неадаптивная система при отклонении параметров с сохранением устойчивости}
  {fig:nonadaptive-stable}

\subsection{Отклонение параметров с потерей устойчивости}

\begin{equation}
  A^{(3)}=\begin{pmatrix}0&1\\\ldots&\ldots\end{pmatrix}
\end{equation}
% Выбрать отклонение параметров, при котором матрица
% A^(3) + b theta_ном^T теряет устойчивость; привести проверку.
\imageplaceholder{nonadaptive_unstable.png}
  {Неадаптивная система при отклонении параметров с потерей устойчивости}
  {fig:nonadaptive-unstable}

\noindent\textbf{Вывод по пункту:} \ldots

\section{Адаптивная система управления}

Настраиваемый регулятор и алгоритм адаптации:
\begin{equation}
  u=\hat{\bm\theta}^{\mathsf T}\bm x+\frac{1}{\kappa}g,
  \qquad
  \dot{\hat{\bm\theta}}=\gamma\bm x\bm b^{\mathsf T}P\bm e,
  \qquad \hat{\bm\theta}(0)=\bm 0,
  \qquad \gamma>0
\end{equation}
Матрица $P=P^{\mathsf T}>0$ определяется из уравнения Ляпунова
\begin{equation}
  A_M^{\mathsf T}P+PA_M=-Q,
  \qquad Q=Q^{\mathsf T}>0
\end{equation}

% Выбрать Q, решить уравнение Ляпунова и проверить положительную
% определённость P. При необходимости привести обоснование закона адаптации.
\begin{equation}
  Q=\ldots,
  \qquad P=\ldots
\end{equation}

Для каждого опыта приводятся три графика: $\bm x(t)$ и $\bm x_M(t)$
на одном поле, ошибка $\bm e(t)$ и параметрическая ошибка
$\tilde{\bm\theta}(t)=\bm\theta-\hat{\bm\theta}(t)$.
% Для изменённого объекта идеальный вектор theta в параметрической ошибке
% определяется из A_M = A^(j) + b theta^T именно для текущей A^(j).

\subsection{Три эксперимента при фиксированном коэффициенте адаптации}

Повторим опыты с матрицами $A^{(1)}$, $A^{(2)}$, $A^{(3)}$ при
\begin{equation}
  \gamma=\gamma_0=\ldots
\end{equation}

\imageplaceholder{adaptive_nominal.png}
  {Адаптивная система при номинальных параметрах объекта}
  {fig:adaptive-nominal}
\imageplaceholder{adaptive_stable.png}
  {Адаптивная система с матрицей объекта $A^{(2)}$}
  {fig:adaptive-stable}
\imageplaceholder{adaptive_unstable.png}
  {Адаптивная система с матрицей объекта $A^{(3)}$}
  {fig:adaptive-unstable}

\subsection{Влияние коэффициента адаптации}

Для номинального объекта сравним три значения коэффициента адаптации:
\begin{equation}
  \gamma_1=\ldots,
  \qquad \gamma_2=\ldots,
  \qquad \gamma_3=\ldots
\end{equation}
% Остальные условия одинаковы. На каждом из трёх графиков сравнить
% результаты всех трёх опытов, подписав соответствующие значения gamma.
\imageplaceholder{adaptive_gamma.png}
  {Сравнение адаптивной системы при трёх значениях $\gamma$}
  {fig:adaptive-gamma}

\subsection{Постоянное задающее воздействие}

% Повторить один из предыдущих адаптивных опытов, заменив только g(t).
\begin{equation}
  g(t)=1,
  \qquad A=\ldots,
  \qquad \gamma=\ldots
\end{equation}
\imageplaceholder{adaptive_constant.png}
  {Адаптивная система при $g(t)=1$}
  {fig:adaptive-constant}

% Оценить ограниченность сигналов, сходимость ошибки слежения и ошибки
% параметров; связать параметрическую сходимость с неисчезающим возбуждением.
% Сравнить скорость настройки при разных gamma и при g(t)=1.
\noindent\textbf{Вывод по пункту:} \ldots

\section{Вывод}

% Обобщить результаты синтеза и моделирования, сравнить неадаптивное
% и адаптивное управление при изменении параметров объекта.
\ldots
